% =====================================================================
%  Introduction to the book
%  (this file is assembled from src/ch00/*.tex by build_chapter.py)
% =====================================================================
\chapter*{Introduction}\label{ch:intro}
\markboth{Introduction}{Introduction}
\addcontentsline{toc}{chapter}{Introduction}

In the spring of 2020, as COVID-19 spread around the world, hundreds of research groups turned to machine
learning. Within months they had published models that diagnosed the disease from chest radiographs and CT
scans, predicted which patients would deteriorate, and forecast the course of the pandemic. Two systematic
reviews examined them afterwards. One of them screened more than two thousand studies on imaging and kept
the 62 of highest quality; it concluded that none of them was of potential clinical use, because of
methodological flaws and biases \citep{roberts2021common}. In some of the data sets, for instance, the
examples without COVID-19 were radiographs of children, so a model could look accurate by recognising
children. The other review found that almost all the published prediction models were at high risk of bias:
small or unrepresentative data, and evaluations that could not be trusted \citep{wynants2020prediction}.

The reviews highlighted serious problems in data selection, study design and evaluation, even when
established learning algorithms were used. The weaknesses lay in everything around the algorithms: how the models were evaluated, on which data they were
tested, whether they treated all patients alike, and whether anyone checked that they would work where
they were going to be used. This book is about everything around the algorithm.

\section*{The life of a model}

A model has a life that starts long before it is trained and continues long after it is deployed
(\cref{fig:lifecycle}). Before it is used, we must know how good it is, and for what purpose
(Unit~\ref{ch:evaluation}); decide whether a combination of models would serve better than any single one
(Unit~\ref{ch:hybrid}); make sure that it is fair, that its decisions can be explained and that it resists
those who want to fool it (Unit~\ref{ch:reliable}); and build it with as much automation as is wise
(Unit~\ref{ch:learning}). Then it must be put in the hands of its users, reliably and repeatably
(Unit~\ref{ch:deployment}), and watched, because the world that produced its training data keeps changing,
until it has to be updated or replaced (Unit~\ref{ch:monitoring}). Recommender systems, customer
relationship management and interactive applications show the whole cycle at work in business
(Unit~\ref{ch:applications}).

\begin{figure}[htbp]
\centering
\begin{tikzpicture}[font=\scriptsize\sffamily,
  st/.style={draw=edmnavy,fill=edmnavy!7,rounded corners=3pt,minimum width=30mm,minimum height=10mm,align=center},
  arr/.style={-{Stealth[length=1.8mm]},edmnavy,thick}]
  \node[st] (u1) at (90:2.6)  {Unit 1\\evaluate};
  \node[st] (u2) at (30:2.9)  {Unit 2\\combine};
  \node[st] (u3) at (-30:2.9) {Unit 3\\make reliable};
  \node[st] (u4) at (-90:2.6) {Unit 4\\automate};
  \node[st] (u5) at (-150:2.9) {Unit 5\\deploy};
  \node[st] (u6) at (150:2.9) {Unit 6\\monitor and update};
  \node[draw=edmbronze,fill=edmgold!20,rounded corners=3pt,minimum width=30mm,minimum height=10mm,align=center]
       (u7) at (0,0) {Unit 7\\applications};
  \draw[arr] (u1) to[bend left=18] (u2);
  \draw[arr] (u2) to[bend left=18] (u3);
  \draw[arr] (u3) to[bend left=18] (u4);
  \draw[arr] (u4) to[bend left=18] (u5);
  \draw[arr] (u5) to[bend left=18] (u6);
  \draw[arr] (u6) to[bend left=18] node[left,font=\tiny\sffamily,align=right] {retrain,\\re-evaluate} (u1);
\end{tikzpicture}
\caption{The life of a model, and the units of this book.}
\label{fig:lifecycle}
\end{figure}

\section*{A first look}

\Cref{lst:firstlook} gives a taste of what is to come with a real data set: the German credit data, 1000
applicants for loans at a German bank, 30\% of whom turned out to be bad credit risks
\citep{hofmann1994statlog}. It trains a logistic regression, the kind of model that takes two lines of code,
and asks it four questions that its accuracy cannot answer.

\begin{pylisting}[label={lst:firstlook}]{Four questions that accuracy does not answer, on real credit data}
import hashlib
import io
from urllib.request import urlopen

import numpy as np
import pandas as pd
from sklearn.compose import make_column_transformer
from sklearn.linear_model import LogisticRegression
from sklearn.model_selection import train_test_split
from sklearn.pipeline import make_pipeline
from sklearn.preprocessing import OneHotEncoder, StandardScaler

# Statlog German credit (UCI), from a mirror pinned to a commit and checked by its hash
URL = ("https://raw.githubusercontent.com/jbrownlee/Datasets/"
       "d20fcb6402ae34e653d4513b00f39257bb37ed7f/german.csv")
SHA256 = "ec12a88b9fc14d74ba646ea0410cf7ff4533bec2eb61652f8ad76796bbfec017"
raw = urlopen(URL).read()
assert hashlib.sha256(raw).hexdigest() == SHA256, "the data file has changed"
data = pd.read_csv(io.BytesIO(raw), header=None)    # 1000 applicants
y = (data[20] == 2).astype(int)             # 1 = the applicant turned out to be bad
X = data.drop(columns=20)
categorical = X.select_dtypes(exclude="number").columns     # coded as A11, A12...
model = make_pipeline(
    make_column_transformer((OneHotEncoder(handle_unknown="ignore"), categorical),
                            remainder=StandardScaler()),
    LogisticRegression(max_iter=1000))
X_tr, X_te, y_tr, y_te = train_test_split(X, y, test_size=0.4, stratify=y,
                                          random_state=1)
p_bad = model.fit(X_tr, y_tr).predict_proba(X_te)[:, 1]

# 1. Is it good? (Unit 1)
print(f"accuracy {np.mean((p_bad >= 0.5) == y_te):.3f}; "
      f"approving everybody: {1 - y_te.mean():.3f}")

# 2. Good for what? Approving a bad applicant costs 5, rejecting a good one 1 (Unit 1)
def cost(threshold):
    reject = p_bad >= threshold
    return (5 * np.sum(~reject & (y_te == 1)) + np.sum(reject & (y_te == 0))) / len(y_te)
print(f"cost per applicant: {cost(0.5):.3f} with threshold 0.5, "
      f"{cost(1 / 6):.3f} with threshold 1/6")

# 3. Good for whom? (Unit 3)
young = (X_te[12] < 25).to_numpy()          # column 12 is the age
approved = p_bad < 1 / 6
print(f"approved: {approved[young].mean():.0%} of the applicants under 25, "
      f"{approved[~young].mean():.0%} of the others")

# 4. Good for how long? If bad applicants become more frequent (Unit 6)
bad_ok = np.mean(p_bad[y_te == 1] < 1 / 6)          # share of bad applicants approved
good_ok = np.mean(p_bad[y_te == 0] < 1 / 6)         # share of good applicants approved
for share in (0.30, 0.50):
    defaults = share * bad_ok / (share * bad_ok + (1 - share) * good_ok)
    print(f"if {share:.0%} of the applicants are bad, {defaults:.1%} of the approved "
          f"loans default")
\end{pylisting}
\begin{pyoutput}
accuracy 0.752; approving everybody: 0.700
cost per applicant: 0.868 with threshold 0.5, 0.530 with threshold 1/6
approved: 29% of the applicants under 25, 43% of the others
if 30% of the applicants are bad, 9.8% of the approved loans default
if 50% of the applicants are bad, 20.1% of the approved loans default
\end{pyoutput}

\begin{itemize}
  \item \emph{Is it good?} Its accuracy is 0.752, which sounds respectable until we notice that approving every
        applicant would score 0.700. Accuracy alone says little (Unit~\ref{ch:evaluation}).
  \item \emph{Good for what?} The data set comes with its own costs: approving a bad applicant costs five times
        more than rejecting a good one. With the default threshold of 0.5 the model costs 0.868 per applicant;
        rejecting every applicant whose probability of being bad exceeds $1/6$, the threshold that these costs
        imply, brings the cost down to 0.530, a reduction of almost 40\% with the \emph{same} model
        (Unit~\ref{ch:evaluation}).
  \item \emph{Good for whom?} The model approves 29\% of the applicants under 25 and 43\% of the others. Is the
        difference justified by their risk, or is the model penalising youth? Answering requires the tools of
        Unit~\ref{ch:reliable}, and in the European Union credit scoring is a \emph{high-risk} application whose
        fairness, documentation and oversight are regulated by law (\cref{sec:aiact}).
  \item \emph{Good for how long?} As an illustrative label-shift scenario, suppose the share of bad applicants rose
        from 30\% to 50\% while the model's approval rates within each class stayed unchanged: the same model would
        approve loans of which 20.1\% default, twice the 9.8\% of today, and the bank would only notice months later,
        when the loans failed to be repaid. Other forms of economic change can alter these rates too. Detecting such changes early is the subject of
        Unit~\ref{ch:monitoring}.
\end{itemize}
None of the four answers is visible in the accuracy, and none of them depends on the learning algorithm.

\section*{Why now}

Three trends make these questions more pressing than ever. First, models now take or prepare decisions about
credit, jobs, health, prices and what we read, watch and buy, often for millions of people at once, so their
mistakes are multiplied. Second, large language models have made it possible for almost anyone to build a
model or an application in an afternoon, and they are themselves models that must be evaluated, deployed and
monitored; we discuss them as tools for automating data science (\cref{sec:llmds}), as components to deploy
(\cref{sec:llmapps}) and as systems to monitor (\cref{sec:llmmonitor}). Third, what used to be good practice is
becoming law: the European Union's AI Act requires those who build high-risk systems to manage their risks,
govern their data, document them, log their operation, keep a human in charge, ensure their accuracy and
robustness, and monitor them once they are on the market (\cref{sec:aiact}). Almost every one of those
obligations corresponds to a unit of this book.

\section*{How to read this book}

Each unit starts with its learning objectives and ends with a summary, bibliographic notes, exercises and
solutions to selected exercises. Boxes highlight key ideas and common pitfalls. The code is in Python, with
scikit-learn and the usual scientific libraries; every listing has been executed, and the output printed below
it is the output it produced. Most examples use synthetic data, generated by the listings themselves, because
synthetic data let us know the truth and change one thing at a time; real data sets appear where they make a
point that synthetic data cannot, as in the listing above. The figures are generated by scripts that import the
same code, so the numbers in the text, in the listings and in the figures agree. All of it is available in the
companion release of the book (see the Preface).

The units can be read in order, as a course, or separately; Unit~\ref{ch:evaluation} is the foundation of the
others, and cross-references point to the sections each unit builds on. The reader is expected to know basic
Python, the elementary concepts of machine learning (training and test data, classification and regression)
and introductory statistics.

